\documentclass[11pt,a4paper]{article}
\usepackage[margin=2.5cm]{geometry}
\usepackage{amsmath,amssymb,bm}
\usepackage{xcolor}
\usepackage[colorlinks=true,linkcolor=blue,urlcolor=blue]{hyperref}
\usepackage{listings}
\usepackage{booktabs}
\usepackage{enumitem}

\lstset{
  basicstyle=\ttfamily\footnotesize,
  breaklines=true,
  frame=single,
  numbers=left,
  numberstyle=\tiny\color{gray},
  keywordstyle=\color{blue},
  commentstyle=\color{green!50!black},
  stringstyle=\color{purple},
}

\title{Julia port of the FSA-v2 model:\\
  closed-loop SMC$^2$ filter + GPU controller}
\author{Coding session writeup}
\date{2026-05-07}

\begin{document}
\maketitle

\section{The Julia FSA-v2 codebase}
\label{sec:codebase}

\subsection{Repository layout}

The model lives at
\path{version_2_Julia/models/fsa_high_res/}. Eight files, mirroring
the Python source under \path{version_2/models/fsa_high_res/}.
\vspace{4pt}

\noindent
\begin{tabular}{@{}lp{0.65\textwidth}@{}}
\toprule
File & Purpose \\
\midrule
\texttt{\_dynamics.jl} & SDE drift + state-dependent diffusion (Stage~G1
                         reparametrisation). \\
\texttt{\_phi\_burst.jl} & Daily-$\Phi \to$ sub-daily morning-loaded
                          burst envelope. \\
\texttt{simulation.jl} & \texttt{INIT\_STATE}, \texttt{DEFAULT\_PARAMS},
                         the four observation samplers, circadian $C(t)$. \\
\texttt{\_plant.jl} & \texttt{StepwisePlant} — \texttt{advance!}
                     (daily $\Phi$ + burst envelope) and
                     \texttt{advance\_subdaily!} (per-bin $\Phi$ direct).
                     \texttt{finalise} writes the psim artefact. \\
\texttt{cpu\_control.jl} & CPU \texttt{build\_control} — RBF + cost
                          functional, debug reference only. \\
\texttt{estimation.jl} & 30 priors + 5 frozen, the locally guided
                          (Pitt--Shephard) \texttt{propagate\_fn},
                          observation log-weights, alignment grid. \\
\texttt{gpu\_pf.jl} & FSA-specific GPU bootstrap-PF kernel +
                     \texttt{FSAGPUTargetBatched}. \\
\texttt{gpu\_control.jl} & FSA-specific GPU cost kernel
                          \texttt{fsa\_cost\_kernel!} +
                          \texttt{FSAControlGPUTarget} +
                          \texttt{make\_log\_density\_fn} closure. \\
\texttt{FSAHighRes.jl} & Module aggregator. \\
\bottomrule
\end{tabular}

\paragraph{Framework additions} (under
\path{julia/SMC2FC/src/}, model-independent):

\begin{itemize}[itemsep=2pt,leftmargin=*]
\item \texttt{Filtering/GPUSegmentedPF.jl} --- shared per-chain weight
  reductions, normalised-weight cumulative sum, systematic resample
  with Liu--West shrinkage, OT sigmoid-blend rescue.
\item \texttt{Control/GPUControlSMC.jl} --- generic parallel-chains
  tempered SMC$^2$, parallel-chains HMC kernel, ChEES-L picker,
  central-FD gradient batcher. Consumes a model-supplied closure
  $\texttt{log\_density\_fn} : \mathbb R^{M\times d} \to \mathbb R^M$.
\end{itemize}

\subsection{Dynamics}

State $\bm X = (B, F, A)$ evolves under the It\^o SDE
\begin{align*}
  \mathrm dB &= \kappa_B \cdot a_B(A) \cdot \Phi \,\mathrm dt
              - \tfrac{B}{\tau_B}\,\mathrm dt + \sigma_B \sqrt{B(1-B)}\,\mathrm dW_B \\
  \mathrm dF &= \kappa_F \Phi\,\mathrm dt - a_F(A)\,\tfrac{F}{\tau_F}\,\mathrm dt
              + \sigma_F \sqrt{F}\,\mathrm dW_F \\
  \mathrm dA &= \mu(B,F)\,A\,\mathrm dt - \eta A^3\,\mathrm dt
              + \sigma_A \sqrt{A}\,\mathrm dW_A,
\end{align*}
with the operating-point reparametrisation $a_B(A)=(1+\varepsilon_A
A)/(1+\varepsilon_A A_{\rm typ})$, $a_F(A)$ the analogue for $\lambda_A$,
and $\mu(B,F) = \mu_0 + \mu_B B - \mu_F F - \mu_{FF}(F-F_{\rm typ})^2$
with $A_{\rm typ}=0.10$, $F_{\rm typ}=0.20$. Truth values use the
effective parameters: $\kappa_B^{\rm eff}=0.01248$, $\tau_F^{\rm eff}=
6.364$, $\mu_F^{\rm eff}=0.26$, $\mu_0^{\rm eff}=0.036$.

\subsection{Plant}

\texttt{StepwisePlant} carries the truth state $(B,F,A)$ and a global bin
counter. \texttt{advance!} expands daily $\Phi$ values via the
morning-loaded Gamma envelope (peak around 10\,am, zero overnight,
integral preserved at $24\cdot \Phi_d$/day) and runs single-step
Euler-Maruyama with the sqrt-It\^o diffusion above.
\texttt{advance\_subdaily!} skips the burst expansion and applies a
caller-supplied per-bin $\Phi$ directly. Both leave the same
\texttt{history} of trajectory + obs that \texttt{finalise} dumps to
disk in psim format.

\subsection{Estimation}

The 30 estimated parameters are declared as a
\texttt{Vector\{Tuple\{Symbol,PriorType\}\}} ordered identically to
Python's \texttt{PARAM\_PRIOR\_CONFIG} so the 30-panel posterior trace
plot lines up column-by-column.
\texttt{propagate\_fn} runs the locally guided (Pitt--Shephard)
proposal: prior predictive $\mathcal N(\mu_{\rm pred}, P_{\rm pred})$
from the G1 drift and state-dependent variance, then sequential scalar
Kalman fusion across the three Gaussian channels (HR, stress,
$\log(\mathrm{steps}+1)$) gated by \texttt{*\_present} masks, then
Cholesky-3 sample. \texttt{obs\_log\_weight\_fn} adds the Bernoulli
sleep log-likelihood. Verbatim port of Python's
\texttt{estimation.py:propagate\_fn}.

\subsection{GPU particle filter}

Model file contains only the FSA-specific propagate kernel — same
locally guided proposal as the CPU version, written as a
KernelAbstractions \texttt{@kernel} that runs in fp32 with one thread
per (chain, particle). Resampling, Liu--West shrinkage and OT rescue
are framework calls (\texttt{Filtering/GPUSegmentedPF.jl}). Public
entry point: \texttt{gpu\_log\_density\_batched(target, U)}, returning
the per-chain marginal log-likelihood.

\subsection{GPU controller}

The control kernel \texttt{fsa\_cost\_kernel!} runs one thread per
(chain $m$, common-random-numbers trial $t$). Per outer bin $k$:
decode $\Phi(t_k) = \Phi_{\max}\sigma\big(c_\Phi + \sum_a
\theta_{m,a}\,\varphi_a(t_k)\big)$ from the (raw, non-normalised)
Gaussian RBF basis; pre-step accumulate the cost integrand
$A_{\rm acc}\mathrel{+}= A\,\Delta t$,
$B_{\rm acc}\mathrel{+}= \max(F-F_{\max},0)^2\,\Delta t$;
then run the substepped EM SDE step under the trial-shared CRN noise.
Cost is Eq~37 of the spec verbatim:
\begin{equation}
  J(\bm\theta_{\rm ctrl}) \;=\; -\!\int_0^T A(t)\,\mathrm dt
  \;+\; \lambda_F \int_0^T \max\bigl(F(t) - F_{\max},\,0\bigr)^2\,\mathrm dt,
  \tag{Eq.\ 37}
\end{equation}
with $\lambda_F=1$, $F_{\max}=0.40$. The model exposes a
\texttt{make\_log\_density\_fn(target)} closure that hands the
per-chain MC-averaged $-J$ to the framework's
\texttt{run\_tempered\_smc\_gpu}.

\subsection{Configuration and wall time}

The reference run --- $T=14$\,d, $h=15$\,min, replan every stride,
$n_{\rm SMC}^{\rm filter}=32$ inner-PF chains $\times K=400$ particles,
$n_{\rm SMC}^{\rm ctrl}=1024 \times n_{\rm inner}=128$ trials, HMC
$\varepsilon=0.2$ / $L_{\rm ChEES}\in\{16,32,64,128,256\}$ /
\texttt{num\_mcmc}=10, target $\beta_{\max}=8$ nats --- completes in
$\sim$780\,s on a single RTX~5090. Python's reference run at the same
configuration takes $\sim$27 minutes.

\section{Bug hunt: what's been ruled out, what remains}
\label{sec:bug-hunt}

The Julia controller settles on a \emph{flat low-$\Phi$} schedule
($\Phi \in [0.10, 0.30]$ across all 27 strides) with mean autonomic
amplitude $\bar A_{\rm MPC} = 0.085$ vs $\bar A_{\rm baseline} = 0.082$
(a $+3.7\%$ gain). Python on the same model reports a
\emph{recovery~$\to$~overload} schedule ($\Phi$ goes
$1.0 \to 0.2 \to 1.2$) with $\bar A_{\rm MPC} = 0.122$
(a $+50\%$ gain). The qualitative shape of the controller's policy
is wrong, not just the magnitude.

This section is the diagnostic timeline. Each hypothesis was tested by
one targeted experiment.

\subsection{Hypothesis 1 (kept): RBF basis was row-normalised}

The original Julia \texttt{build\_rbf} row-normalised the design matrix
into a partition of unity. Python's
\texttt{smc2fc/control/rbf\_schedules.py:design\_matrix} returns the
\emph{raw} Gaussian basis (no normalisation). With row-normalisation
the decoded $\Phi(t)$ collapses to a weighted average of $\theta_a$
values, smoothing out any per-anchor variation regardless of the
gradient signal.

\textbf{Test:} removed normalisation in both
\path{cpu_control.jl:build_rbf} and the GPU target's RBF construction.
\textbf{Outcome:} $\Phi$ moved from a flat $\sim 0.45$ to a flat
$\sim 0.20$, matching the spec's text description of a
``rest-leaning $\Phi \in [0.25, 0.31]$''. Real fix, kept in code.
\emph{It did not produce the recovery~$\to$~overload pattern, but it
was a real bug.}

\subsection{Hypothesis 2 (kept): three small kernel mismatches vs Python}

Found while reading the Julia kernel against Python's
\texttt{control.py:cost\_fn} line-by-line:

\begin{enumerate}[itemsep=2pt]
\item A NaN guard at the end of every bin reset the state to
  $(0.05, 0.30, 0.10)$ if any of $B,F,A$ became non-finite.
  Removed --- the boundary reflection already handles the legitimate
  cases, and the guard was masking real numerical issues.
\item The cost diffusion used $\sqrt{A + \varepsilon_A}$
  ($\varepsilon_A=10^{-4}$). Python's controller path uses
  $\sqrt{\max(A, 0)}$ (the $\varepsilon_A$ form lives only in the
  plant's \texttt{noise\_scale\_fn}). Aligned.
\item The cost integrand $A_{\rm acc}$ was being accumulated AFTER the
  EM step. Python uses the PRE-step state. Switched.
\end{enumerate}

\textbf{Outcome:} small numerical changes, no qualitative effect on the
controller's optimum. Kept anyway --- they bring the kernel into bit-for-bit
correspondence with Python's algorithm.

\subsection{Hypothesis 3 (discarded): bursty / smooth $\Phi$ asymmetry between MPC and baseline plants}

The bench was calling \texttt{advance!} (which expands daily $\Phi$
through the burst envelope) for the baseline plant and
\texttt{advance\_subdaily!} (which takes per-bin $\Phi$ directly) for
the MPC plant. Under constant $\Phi=1$, a Monte-Carlo plant test gave
$\bar A_{\rm bursty}=0.082$ vs $\bar A_{\rm smooth}=0.100$ — a 22\,\%
gap that suggested the MPC trajectory was being measured on a different
fast-dynamics regime from the baseline.

\textbf{Test:} aggregated the controller's per-bin $\Phi$ to daily
values and routed both plants through \texttt{advance!} (so both run
on the bursty envelope). \textbf{Outcome:} the bench produced
essentially identical numbers ($\bar A_{\rm MPC}=0.085$, $\Phi$ still
flat at $\sim 0.20$). The baseline-plant asymmetry was real and worth
fixing, but it was \emph{not} the source of the controller's wrong
optimum. \emph{Hypothesis discarded.}

\subsection{Hypothesis 4 (discarded): bug in the framework controller}

Suspicion: the parallel-chains tempered SMC$^2$ + ChEES-HMC machinery
in \path{julia/SMC2FC/src/Control/GPUControlSMC.jl} might have a
gradient-following bug or a tempering-bisection bug that prevents it
from finding off-prior-mean optima.

\textbf{Test:} ran the framework controller against three synthetic
analytic costs with \emph{known} optima. Pure CPU log-density closures
($\mathbb R^{M\times d} \to \mathbb R^M$), no model and no GPU kernel.
Driver: \path{version_2_Julia/tools/test_framework_controller.jl}.

\begin{itemize}[itemsep=4pt,leftmargin=*]
\item \textbf{Quadratic bowl} $J(\bm\theta)=\tfrac12 \|(\bm\theta -
  \bm\theta^\star)/\sigma_J\|^2$ with $\bm\theta^\star \neq \mathbf 0$.
  Framework returned posterior mean $0.755\cdot \bm\theta^\star$,
  matching the analytic Bayesian-shrinkage prediction at the
  auto-calibrated $\beta_{\max}=1.37$ within rounding. \emph{Pass.}
\item \textbf{Anti-symmetric reward} $J(\bm\theta) = -\bm w^\top
  \bm\theta + \gamma\|\bm\theta\|^2$ with $\bm w$ anti-symmetric across
  the eight RBF anchors (negative early, positive late) — the same
  spatial structure as the FSA recovery~$\to$~overload schedule.
  Framework returned a posterior mean with the correct anti-symmetric
  signs and ratios, shrunk by $\sim 25\%$ — same Bayesian-shrinkage
  signature. \emph{Pass: the framework finds spatially structured
  optima when they exist.}
\item \textbf{Bimodal cost} with one ``flat'' mode at
  $\bm\theta_a = -\mathbf 1$ and one ``recovery~$\to$~overload'' mode
  at $\bm\theta_b = (-1.5,-1.5,-1.5,0,0,+1.5,+1.5,+1.5)$, with
  $\bm\theta_b$ the global mode. Tested three initialisations: from
  prior, from $\bm\theta_a$, from $\bm\theta_b$. All three converged to
  the same prior-weighted compromise — not a defect, but the consequence
  of $\beta_{\max}$ auto-calibrating to $0.367$ under the wide prior
  (the prior dominates the tempered posterior at small $\beta$).
\end{itemize}

\textbf{Outcome:} the framework's controller logic is correct; in
particular it can find anti-symmetric optima when they are present
(Test~2). \emph{Hypothesis discarded.}

\subsection{Hypothesis 5 (discarded): HMC stuck in the wrong basin on FSA}

Even if the cost surface contained both a flat-low and a
recovery~$\to$~overload basin and the latter were deeper, an
HMC chain initialised in the larger flat-low basin might never escape.

\textbf{Test:} ran the FSA bench with the controller's HMC initialised
at a constructed recovery~$\to$~overload $\bm\theta_{\rm seed}$ (per-anchor
target $\Phi = (1.0, 0.2, 0.2, 0.4, 0.7, 0.9, 1.1, 1.2)$, decoded
through the inverse RBF formula, $\sigma_{\rm jitter}=0.3$ around the
seed). Same controller config, same wall ($\sim 740$\,s).

\textbf{Outcome:} every replan still settled at $\bar\Phi \in [0.27,
0.30]$, identical final state $\hat x = (0.064, 0.071, 0.096)$ to the
cold-start runs. HMC initialised at the recovery~$\to$~overload point
\emph{moved away} from it and converged to flat-low. The
recovery~$\to$~overload region is not a local minimum at all in the
Julia FSA cost surface --- the gradient there points toward flat-low.
\emph{Hypothesis discarded.}

\subsection{Hypothesis 6 (kernel + decoder OK on monotone reward)}

A trivial-cost end-to-end test that isolates the FSA SDE + GPU kernel +
RBF decoder + framework HMC, removing all the Stuart--Landau and
F-coupling nonlinearities of the original cost.

\paragraph{Test cost.} Replace Eq~37 by
\[
J(\bm\Phi) \;=\; -\!\int_0^T B(t)\,\mathrm dt,\qquad \lambda_F = 0
\quad(\text{no fatigue penalty}).
\]
Under the FSA SDE, $\mathrm dB/\mathrm dt = \kappa_B \cdot a_{B}(A) \cdot \Phi
- B/\tau_B$ is monotone in $\Phi$. The optimum is to pin $\Phi=\Phi_{\max}=3$
for every bin, driving $B$ toward its highest equilibrium
$\kappa_B\,\Phi_{\max}\,\tau_B$.

\paragraph{Implementation.} The cost kernel is parametrised with
state-integrand weights $(w_B, w_F, w_A)$ (default $(0,0,1)$, i.e.\ Eq~37).
Setting $(1, 0, 0)$ and $\lambda_F = 0$ gives the trivial cost.
Driver: \path{version_2_Julia/tools/test_max_B_only.jl}.

\paragraph{Result.} The framework returns posterior-mean
$\bar{\bm\theta} = (+1.57, +2.12, +1.92, +1.77, +1.71, +1.52, +1.04, +0.42)$
which decodes to a per-bin $\Phi$ schedule with:
\begin{itemize}[itemsep=1pt]
\item day 0 to day 8: $\Phi(t) \in [2.68, 2.92]$ — essentially
  pinned to $\Phi_{\max}$.
\item day 10 to day 14: $\Phi$ tapers down to $1.79$ — Bayesian
  shrinkage near the boundary anchors where the prior pull is more
  visible.
\item summary: $\Phi_{\rm mean} = 2.73$ vs $\Phi_{\max} = 3.0$.
\end{itemize}
Six tempering levels, auto-calibrated $\beta_{\max} = 12.5$
(large because $\int B$ has small variance under the prior cloud, so
cost-std is small), 55\,s wall.

\paragraph{Conclusion.} When the optimum is monotone in $\Phi$, the
controller finds it. End-to-end this rules out:
\begin{itemize}[itemsep=1pt]
\item Bug in the FSA SDE integration ($\mathrm dB/\mathrm dt$ is correct).
\item Bug in the RBF decoder ($\theta_a > 0$ correctly produces
  high $\Phi$ at anchor $a$).
\item Bug in the GPU-kernel cost-accumulation pathway
  ($\int B\,\mathrm dt$ is correctly assembled).
\item Bug in the framework parallel-chains HMC.
\end{itemize}
The remaining root cause is therefore confined to the A-cost surface
specifically --- the Stuart--Landau cubic $-\eta A^3$, the
$F$-coupling through $\mu(B, F)$, or some numerical interaction between
them that displaces the global optimum from recovery~$\to$~overload to
flat-low in Julia versus Python.

\subsection{Hypothesis 7 (root cause confirmed): fp32 cancellation in the G1 drift}

Triggered by the observation that the v1 reference controller
(\path{version_1_Julia/models/fsa_high_res/gpu_control.jl}, working
end-to-end on the same SDE structure) uses the \emph{pre-G1} drift
form, while my v2 uses the G1-reparametrised form. The two are
mathematically equivalent at the truth parameters --- v1 mu and v2 mu
both evaluate to $-0.031$ at $(B,F,A)=(0.05, 0.30, 0.10)$. The G1
reparametrisation cannot be unwound on the v2 side because v2 needs
it for FIM-rank identifiability of the static parameters in the
\emph{filter}. So the controller has to use the v2 form.

The v2 form contains arithmetic that is more cancellation-prone in
fp32 than the v1 form:
\begin{itemize}[itemsep=2pt]
\item $(F - F_{\rm typ})^2$ subtracts magnitudes of order $F_{\rm typ}=0.20$
  then squares the small remainder, losing fp32 mantissa bits.
\item $(1 + \varepsilon_A A) / (1 + \varepsilon_A A_{\rm typ})$ divides
  near-equal numbers when $A \approx A_{\rm typ}$. Both numerator and
  denominator are $\approx 1.04$; the ratio is determined by their
  small difference. Same for the $\lambda_A$ term in $a_F(A)$.
\item Larger individual terms in $\mu$ that have to cancel:
  $\mu_F^{\rm eff} F = 0.078$ at $F=0.30$ vs $\mu_F^{\rm v1} F = 0.030$.
\end{itemize}
Each substep evaluation of the drift accumulates rounding from these
sources. Over the controller's cost rollout
($n_{\rm steps} \times n_{\rm substeps} = 1344 \times 4 = 5376$
evaluations per trial, $n_{\rm inner}=128$ trials per chain), the
biased rounding shifts integrated A by enough to displace the global
cost minimum from the recovery$\to$overload schedule to flat-low.

\paragraph{Test.} The drift block inside the substep loop of
\path{gpu_control.jl:fsa_cost_kernel!} was promoted to fp64 (state and
accumulators stay fp32; only the per-substep arithmetic that produces
\texttt{drift\_B}, \texttt{drift\_F}, \texttt{drift\_A} runs in fp64).
A single planning call at the canonical init state
$(B,F,A)=(0.05,0.30,0.10)$ over the full $T=14$\,d horizon was run,
with a light controller config
($n_{\rm SMC}=64$, $n_{\rm inner}=16$, $\texttt{num\_mcmc}=3$,
$L_{\rm leap}=8$) to keep wall time tractable on consumer fp32:fp64
silicon (RTX~5090's fp64 throughput is $\sim 1/64$ of fp32). Driver:
\path{version_2_Julia/tools/test_one_plan_fp64.jl}.

\paragraph{Result.} Posterior-mean $\bar{\bm\theta}$ decoded to:

\begin{tabular}{@{}rl@{\qquad}rl@{}}
day 0 & $\Phi = 0.23$ & day 8  & $\Phi = 0.31$ \\
day 1 & $\Phi = 0.13$ & day 9  & $\Phi = 0.36$ \\
day 2 & $\Phi = 0.10$ & day 10 & $\Phi = 0.42$ \\
day 3 & $\Phi = 0.10$ & day 11 & $\Phi = 0.52$ \\
day 4 & $\Phi = 0.12$ & day 12 & $\Phi = 0.67$ \\
day 5 & $\Phi = 0.15$ & day 13 & $\Phi = 0.84$ \\
day 6 & $\Phi = 0.20$ & day 14 & $\Phi = 0.98$ \\
day 7 & $\Phi = 0.25$ & & \\
\end{tabular}
\vspace{2pt}

\noindent
This is the \emph{recovery~$\to$~overload} pattern. $\Phi$ drops to
its floor in days 1--3 (clearing the high initial fatigue), holds low
through day 5--6, then ramps monotonically up to nearly $\Phi_{\max}/3$
by day 14. Compare to the fp32 production run, which gave a flat
$\Phi \approx 0.20$ across all 14 days.

The auto-calibrated $\beta_{\max}$ also moved dramatically: 117 in fp64
vs $\sim$12 in fp32, indicating the cost has $\sim 10\times$ larger
variance across the prior cloud once the cancellation is removed ---
the cost surface is genuinely sharp, the fp32 noise was washing it out.

\paragraph{Hypothesis confirmed.} The flat-low controller behaviour
was caused by accumulated fp32 cancellation in the v2 G1 drift.

\subsection{Outstanding: production-grade fix}

The fp64-promotion test confirms the diagnosis but is not a
production fix --- it costs $\sim 30\times$ more wall-time on the 5090
(consumer Blackwell's fp64 throughput is $1/64$ of fp32). The full
bench at production config ($n_{\rm SMC}=1024, n_{\rm inner}=128$)
would take $\sim 15$ hours under fp64 substep arithmetic.

The proper fix is an algebraic reformulation of the v2 G1 drift that
avoids the cancellation hot-spots while staying mathematically
identical:

\begin{itemize}[itemsep=2pt]
\item Rewrite $\mu = \mu_0 + \mu_B B - \mu_F F - \mu_{FF}(F-F_{\rm typ})^2$
  as a Horner-style polynomial in $F$ that fuses the constant terms:
  $\mu = (\mu_0 - \mu_{FF} F_{\rm typ}^2) + \mu_B B + (2\mu_{FF}F_{\rm typ} - \mu_F) F - \mu_{FF} F^2$.
  This is identically the v1 polynomial in $F$ once the constant and
  linear coefficients are pre-combined --- an algebraic rearrangement
  that the JIT may or may not perform on its own.
\item Pre-compute $a_{\rm typ}^{-1} \equiv 1/(1+\varepsilon_A A_{\rm typ})$
  once at target construction, then evaluate $a_B(A) = (1+\varepsilon_A A) \cdot a_{\rm typ}^{-1}$
  --- one multiply, no division, no near-equal subtraction.
\item Same trick for $a_F(A)$.
\end{itemize}

After the reformulation, the production fp32 kernel should reach the
same $\beta_{\max}$ and the same posterior-mean schedule as the fp64
diagnostic, at the original $\sim 800$\,s wall time.

\paragraph{Update 2026-05-07 (post-application).}
The Horner reformulation and the precomputed $a_{\rm typ}^{-1}$
factors were applied to \path{gpu_control.jl:fsa_cost_kernel!} (commit
in working tree). Three verifications followed:
\begin{itemize}[itemsep=2pt]
\item Bit-for-bit at $\theta=0$ and the hand-crafted
  $\theta_{\rm RO}$ — Julia GPU-fp32 and Python fp64 still agree to a
  max relative difference of $3\times 10^{-6}$.
\item Light-config open-loop ($n_{\rm SMC}{=}64$, $n_{\rm inner}{=}16$):
  the controller still finds the recovery$\to$overload optimum,
  $\Phi$ ramps $0.10 \to 0.95$ over $14\,$d, $\sim$15\,s wall.
\item \textbf{Heavy-config closed-loop ($n_{\rm SMC}{=}1024$,
  $n_{\rm inner}{=}128$, 27 replans): the bench still settles to
  flat-low.} Per-replan plan means $\bar\Phi_{\rm plan} \in
  [0.26, 0.39]$ (recovery$\to$overload-shaped \emph{plans}), but
  cumulative applied $\Phi$ stays in $[0.10, 0.35]$ and final
  $\bar A_{\rm MPC}=0.085$ vs baseline $0.082$ --- the same flat-low
  numbers as before \S2.8. $701$\,s wall time, matching the reference.
\end{itemize}

\textbf{Conclusion: \S2.8 alone is necessary but not sufficient.}
The cancellation hot-spots are gone (which is why the light-config
open-loop test continues to recover the right optimum), but the heavy
closed-loop is still broken. See \S2.10 for the standing hypothesis and
the staged plan that supersedes the single \S2.8 fix.

\subsection{Summary of what was ruled out, and what fixed it}

\begin{tabular}{@{}p{0.16\textwidth}p{0.5\textwidth}p{0.27\textwidth}@{}}
\toprule
Hypothesis & Test & Outcome \\
\midrule
1. RBF row-norm    & Removed, re-ran bench         & Real bug. Kept.    \\
2. Kernel mismatches (NaN guard, $\varepsilon_A$, POST-step) & Aligned to Python, re-ran & Cosmetic. Kept. \\
3. Bursty/smooth $\Phi$ asymmetry & Routed MPC plant through \texttt{advance!} & No effect. Discarded. \\
4. Framework controller bugs & 3 synthetic costs with known optima & Framework correct. \\
5. HMC stuck in wrong basin & Seeded HMC at recovery$\to$overload  & HMC drifts away. Discarded. \\
6. SDE / decoder / kernel infrastructure & Max-$\int B$ trivial reward & $\Phi \to \Phi_{\max}$ correctly. \\
7. fp32 cancellation in G1 drift & fp64-promoted substep arithmetic  & Recovery$\to$overload restored at light config. \\
8. \S2.8 algebraic rearrangement (Horner $\mu$ + $a_{\rm typ}^{-1}$) & Applied; bit-for-bit + light open-loop + heavy closed-loop & Light open-loop OK; heavy closed-loop still flat-low. \S2.8 was necessary but not sufficient. \\
9. Bench was using shrinking remaining-bench horizon as planning horizon & Replaced with fixed $H_{\rm plan} = T_{\rm bench}$ at every replan (mirrors Python) & Real bug. Applied $\Phi$ now has Python's qualitative U-shape (1.0 day 0, 0.2 mid, 0.78 by day 13). \emph{But} amplitude is too low: $\bar A_{\rm MPC}$ moved only from 0.085 to 0.087 vs Python's 0.122. Necessary, not sufficient. \\
10. Why does Julia under-shoot Python's $\bar A$ even with right shape? & Apply Python's exact daily-$\Phi$ schedule to Julia's plant (\path{tools/test_apply_python_schedule.jl}). & \textbf{Plant-design difference, by design.} Julia uses one 24-bin daily envelope; Python uses two 12-bin per-stride envelopes. Same daily-$\Phi$ → different sub-day $F$ trajectories → Julia's $\bar A = 0.087$ vs Python's 0.122 under \emph{identical} schedules. Accepted: Julia's plant is the intended design. \\
\bottomrule
\end{tabular}
\vspace{4pt}

The cancellation hot-spots in the v2 G1 drift's
$(F-F_{\rm typ})^2$ and $(1+\varepsilon_A A)/(1+\varepsilon_A A_{\rm typ})$
are gone after \S2.8. The light-config open-loop test confirms the
cost-surface shape is right at $n_{\rm inner}=16$. But at production
$n_{\rm inner}=128$ the heavy closed-loop bench still settles to
flat-low, so a residual fp32 bias must survive elsewhere in the
kernel. \S2.10 lays out the staged-promotion plan for finding the
minimum sufficient fp64 promotion that fixes the heavy closed-loop.

\subsection{Hypothesis 8 (in progress): residual fp32 bias surviving \S2.8 --- staged fp64 promotion}

\textbf{Standing hypothesis after \S2.8 was applied and didn't fix the
closed-loop:} the Horner reformulation removed the worst (cancellation-
prone) fp32 ops, but smaller systematic biases survive in the still-fp32
parts of the kernel. The mechanism that's consistent with all the
observed numbers:

\begin{itemize}[itemsep=2pt]
\item At \textbf{light config} ($n_{\rm inner}=16$) the per-cost MC
  noise is $\sim\sigma/\sqrt{16}$. This noise dominates whatever
  systematic fp32 bias is left, so the cost surface seen by HMC is
  approximately the same as Python's at low signal-to-noise --- the
  global recovery$\to$overload basin still wins.
\item At \textbf{heavy config} ($n_{\rm inner}=128$) the MC noise
  shrinks by $\sqrt{8}\approx 2.8\times$. The systematic fp32 bias
  that was buried in noise becomes the dominant distortion, and
  because it's plan-dependent it warps the cost surface enough to
  displace the global optimum.
\end{itemize}

This is consistent with: bit-for-bit at probe $\theta$ values (where
the cost is point-evaluated and noise/bias both irrelevant) matches
Python; light-config open-loop matches Python; heavy-config closed-loop
diverges.

\paragraph{Three candidate locations for the residual fp32 bias.}
Inside \path{gpu_control.jl:fsa_cost_kernel!} after \S2.8, the
remaining fp32 paths are:

\begin{enumerate}[itemsep=2pt,leftmargin=*]
\item \textbf{Cost accumulators} (\texttt{A\_acc}, \texttt{barrier\_acc}).
  Sum of $1344$ fp32 increments to a total reaching $\sim 1.4$. By the
  tail of the sum the relative add ratio
  $|A\,\Delta t|/|A_{\rm acc}| \approx 10^{-3}$ --- below the fp32
  mantissa floor for the smallest dropped bits. The bias is
  \emph{plan-dependent} (plans driving $A$ higher accumulate more
  bits of loss), so it distorts the cost \emph{surface}, not just the
  level. \textbf{Cheapest to fix} --- two scalar promotions per thread,
  no per-substep cost.
\item \textbf{State accumulation}, $B \mathrel{+}= \Delta t \cdot
  \dot B$ etc., repeated $5376$ substeps. Same compounding-rounding
  pattern, on the state itself. If state drifts off, downstream $A$
  diverges.
\item \textbf{$\mu_{\rm bif}$ Horner sum} --- four signed terms of
  similar magnitude added in fp32 each substep. Less likely to be
  plan-dependent at the magnitudes involved (all four terms
  $\mathcal O(0.01\text{--}0.1)$) but worth checking if (1)+(2) don't
  suffice.
\end{enumerate}

\paragraph{Staged methodology.}
Promote one location at a time, measure each step against three
verifications:

\begin{enumerate}[itemsep=2pt,leftmargin=*,label=\textbf{Step \arabic*.}]
\setcounter{enumi}{0}
\item Restore the \S2.8 production baseline (fp32 + Horner) and
  re-confirm the three known outcomes: bit-for-bit at probe
  $\theta$ values $\le 3\times 10^{-6}$; light open-loop
  recovery$\to$overload; heavy closed-loop flat-low.
\item Promote \textbf{cost accumulators} to fp64 only.
  Verifications: bit-for-bit (should improve slightly), light
  open-loop (must remain recovery$\to$overload), heavy closed-loop
  (the diagnostic --- does $\Phi$ shift?).
\item If Step 1 didn't shift the heavy closed-loop, additionally
  promote the \textbf{state} ($B, F, A$) to fp64. Same verification
  suite.
\item If Step 2 didn't shift it either, additionally promote the
  \textbf{drift computation} ($\mu_{\rm bif}$, $a_{\rm factor}$, the
  drift outputs) to fp64. This is the upper bound; the failed
  full-fp64 conversion already showed open-loop matches Python at
  this point, so heavy closed-loop must too --- otherwise the
  framing of ``residual fp32 bias displaces the optimum'' was wrong
  and the bug is somewhere orthogonal (orchestration, plant
  $\leftrightarrow$ kernel mismatch, etc.).
\item Pick the cheapest set that produces recovery$\to$overload in
  the heavy closed-loop and lock it in with a comment block citing
  the empirical comparison.
\end{enumerate}

Stop conditions: stop early if Step 1 already fixes the heavy bench
(ship it); stop and reassess if any step regresses the open-loop test
(means the change broke something orthogonal); stop after Step 3
regardless and shift the investigation back to orchestration.

\paragraph{Verification target.} The Julia heavy closed-loop trace
plot \path{E5_full_mpc_T14d_traces.png} should qualitatively match
the Python reference at
\path{version_2/outputs/fsa_high_res/g4_runs/T14d_replanK2_h60min_no_infoaware/E5_full_mpc_T14d_traces.png}:
applied $\Phi$ visibly U-shaped (high day-0, low days 1--5, ramping
back up to $\sim 1.0$ by day 13), $F$ decaying then turning back up at
the overload phase, $\bar A_{\rm MPC} \gtrsim 0.115$ (Python target
$0.122$), F-violation $\le 1\%$.

The corresponding implementation plan is
\path{claude_plans/Localising_residual_fp32_bias_in_the_v2_GPU_cost_kernel_2026-05-07_1956.md}.

\subsection{Hypothesis 9 (real bug, partial fix): receding-horizon MPC was using the wrong horizon}

\textbf{A real MPC formulation bug, fixed --- but not the full root cause.} The bench's per-replan plan was being
constructed over the \emph{shrinking remaining-bench horizon}, not over
a \emph{fixed planning horizon}. These are two distinct concepts that
the original code conflated:

\begin{itemize}[itemsep=2pt]
\item $T_{\rm bench}$ --- the closed-loop \emph{bench duration} (= 14 d
  here). How long the experiment runs.
\item $H_{\rm plan}$ --- the controller's \emph{planning horizon}. How
  far ahead the controller looks at each replan.  Should be HELD FIXED
  throughout the bench at a value chosen for the dynamics (here $H_{\rm
  plan} = \tau_B = 14\,$d, one chronic-B time constant --- enough
  lookahead for the recovery$\to$overload trade-off to be visible).
\end{itemize}

The original Julia bench tied $H_{\rm plan} = T_{\rm bench} - t_{\rm
elapsed}$. By stride 20 of a 14-day bench (= day 10), $H_{\rm plan}$
had shrunk to 4 days. The cost-surface optimum at the current state
(low $F$, moderate $B$) over a 4-day window collapses to ``recover''
--- there's no time left for the controller to commit to a final
overload phase that would actually pay back in $\int A\,\mathrm{d}t$.
Plan day-0 stays at the recovery start ($\Phi \approx 0.20$); the
bench commits this and replans the next stride from a state still in
the recovery basin; same plan shape; same $\Phi \approx 0.20$ committed;
$\dots$ for the rest of the bench. Cumulative applied $\Phi$ stuck in
[0.10, 0.30].

Python's \path{bench_smc_full_mpc_fsa.py} line 389 hard-codes
\path{plan_horizon_days=T_total_days} at every replan: a FIXED 14-day
lookahead, regardless of how much bench remains. At late strides the
plan extends \emph{past} the bench's end; only the first stride is
committed before the next replan, so the past-the-end portion is just
discarded.  With a long-enough lookahead, the cost optimum from the
late-bench state (low $F$, moderate $B$) is the U-shape: overload
immediately (because $F$ is low and there's room to push $B$ further),
brief recovery, then a second overload. Plan day-0 = high; applied
$\Phi$ rises to $\sim 1.2$ at the bench's end.

\paragraph{Verification by isolated open-loop test.}
A 14-day open-loop plan from the late-bench state $(B, F, A) =
(0.075, 0.092, 0.090)$ with the heavy controller config returns
posterior-mean
$\bar\theta = (+0.40, -0.08, -0.66, -0.87, -0.95, -0.62, -0.42, -0.17)$
which decodes to $\Phi$ at days $\{0, 2, 4, 6, 8, 10, 12, 14\} =
\{1.21, 0.82, 0.39, 0.24, 0.23, 0.33, 0.50, 0.73\}$ --- the
\emph{U-shape} that matches Python's reference plot. Plan day-0
$= 1.21$. From the same state, a 4-day open-loop plan gives plan
day-0 $\approx 0.22$ (recovery). The cost optimum's plan day-0 is
plan-horizon-dependent in exactly the way the failure mode predicts.

\paragraph{The partial fix.}
\path{version_2_Julia/tools/bench_smc_full_mpc_fsa_gpu.jl} now replans
with \texttt{ctrl\_n\_steps = round(args["T-days"] / DT\_DAYS)} at
every replan boundary, regardless of \texttt{mpc\_plant.t\_bin}. The
plan covers a fixed lookahead; only the first stride is committed.
Combined with Hypothesis 8's end-of-stride replan strategy (initial
baseline $\Phi=1.0$ for the first $K$ strides, then replan every $K$),
the closed-loop bench's applied $\Phi$ now has the qualitative
U-shape that matches Python's reference --- $\Phi = 1.0$ on day 0,
drop to $\sim 0.2$ during recovery (days 2--5), ramp up to $\sim 0.78$
by day 13.

\paragraph{What's still wrong.}
\textbf{The applied $\Phi$ shape is right but its amplitude is too low,
and the $A$ trajectory does not climb to Python's level.}

\begin{itemize}[itemsep=2pt]
\item Python's reference: applied $\Phi$ peaks at 1.2 by day 13,
  $A(\mathrm{MPC})$ peaks at 0.17, $\bar A_{\rm MPC} = 0.122$ vs
  baseline $0.081$ ($+51\%$ gain).
\item Julia after this fix: applied $\Phi$ peaks at 0.78 by day 13,
  $A(\mathrm{MPC})$ stays in $[0.08, 0.10]$ and actually \emph{drops}
  during recovery, $\bar A_{\rm MPC} = 0.087$ vs baseline $0.085$
  (a +2.3\% gain).
\end{itemize}

So the controller is shaping the schedule correctly (recovery
followed by overload) but is committing too little overload --- not
enough late-phase $\Phi$ to drive $B$ high enough to lift $\mu(B,F)$
into the regime where $A^* = \sqrt{\mu/\eta}$ exceeds the initial
$A = 0.10$. There is at least one more bug between the controller
and the plant; \S2.11 picks up the investigation.

\paragraph{What was ruled out by getting here.}
\S2.1 -- \S2.7 had narrowed the bug to ``somewhere in the cost
surface near the $A$-cost integrand under the v2 G1 drift'', and
\S2.8 / \S2.10 tried algebraic and staged-fp64 fixes. None of those
touched the closed-loop optimum. \S2.9 (this section) is a real bug
--- the receding-horizon was using the wrong horizon, and the fix
materially shifts the closed-loop schedule into the right qualitative
shape --- but it does not fully reproduce Python's gain. So the
kernel, the framework, the precision setup, and \S2.8 are all
correct (open-loop tests confirm); \S2.9's MPC formulation is also
now correct; but at least one more bug is still affecting the
closed-loop trajectory.

\subsection{Hypothesis 10 (root cause): bench was feeding the plant a re-bursted daily-average of the controller's smooth schedule}

\textbf{The actual bug}, found by side-by-side reading against the
ground-truth driver
\path{tools/test_max_A_with_F_barrier.jl}: the bench was taking the
controller's per-bin $\Phi$ slice for each stride, \emph{averaging it
to a single daily value}, and then \emph{re-expanding it through the
morning-loaded burst envelope} before passing to the plant.

\begin{verbatim}
# wrong (bench, before fix):
daily_phi_for_stride = mean(phi_stride)     # destroy controller's shape
advance!(mpc_plant, advance_bins, [daily_phi_for_stride])  # re-burst
\end{verbatim}

That replaces the controller's actual schedule with a different
signal: a daily-averaged $\Phi$ refracted through a Gamma-shaped
morning burst. The plant then runs an SDE under that re-bursted
signal, not the smooth schedule the controller's cost kernel was
optimising against. Pure plant $\leftrightarrow$ controller mismatch.

The ground-truth driver (\S 2.10 above) calls
\path{advance_subdaily!} with the controller's per-bin $\Phi$
directly:

\begin{verbatim}
# ground truth (correct):
advance_subdaily!(mpc_plant, Float32.(slice))  # controller's exact phi(t)
\end{verbatim}

Same call now in the bench:

\begin{verbatim}
# bench, after fix:
advance_subdaily!(mpc_plant, Float32.(phi_stride))
advance_subdaily!(base_plant, fill(Float32(1.0), advance_bins))
\end{verbatim}

After this fix the bench's mean $A_{\rm MPC}$ is essentially the
same as the ground-truth open-loop driver's ($0.087$ vs $0.091$,
within MC noise --- the ground truth has cleaner numbers because
it commits one plan rather than re-optimising every $K$ strides).
The qualitative $\Phi$ trace now has the U-shape Python's reference
shows (baseline at days 0--1, recovery to $\Phi \approx 0.2$ days
2--5, ramp with peaks reaching $\Phi \approx 1.4$ by day 10).

\paragraph{Why this is the actual root cause.}
The diagnostic chain through \S 2.1--\S 2.9 was repeatedly hitting a
``flat-low applied $\Phi$'' symptom. Every time we fixed something
upstream of the plant (RBF norm, kernel mismatches, fp32 cancellation,
algebraic Horner form, planning horizon, replan strategy) the
applied schedule got closer to recovery$\to$overload but $A$
never grew. With the plant call corrected, the
controller's smooth $\Phi(t)$ now actually drives the plant. The
A/B test on burst-envelope modes (\path{--plant single_daily} vs
\path{per_stride}) gave the same mean $A = 0.087$ in both modes
\emph{because both modes are wrong} relative to what the controller
expects --- the controller wants a smooth signal, not any flavour
of bursty refraction. The fix is to drop the bursty refraction
entirely and pass the per-bin $\Phi$ as-is.

\paragraph{Status.}
\textbf{This was the residual bug after \S 2.9.} The full bug list:

\begin{enumerate}[itemsep=2pt,leftmargin=*]
\item \S 2.1 RBF row-normalisation (kept fix --- real bug).
\item \S 2.2 NaN guard / $\varepsilon_A$ / POST-step accumulation
  (kept fixes --- cosmetic).
\item \S 2.9 Receding-horizon used shrinking remaining-bench
  horizon (kept fix --- real bug).
\item \S 2.10 (this section) Bench was re-bursting the controller's
  smooth schedule before passing to the plant (kept fix --- real
  bug, the residual bottleneck).
\end{enumerate}

\S 2.3, \S 2.4, \S 2.5, \S 2.6 were ruled-out hypotheses (controller
machinery is fine). \S 2.7, \S 2.8 / \S 2.11 chased an fp32
cancellation hypothesis that the staged-promotion test (Step 1
fp64 accumulators, Step 2-alt Kahan + FMA) ruled out --- the
algebraic-stable kernel is kept anyway as a free correctness win,
but it was not load-bearing for the optimum.

\paragraph{Open-loop mode in the bench.}
After fixing the plant call, the closed-loop bench's applied $\Phi$
still oscillates because each replan re-finds a slightly different
optimum, and the cost surface has multiple plausible basins HMC can
land in. To verify the bench is now structurally equivalent to the
ground-truth driver, the bench has an \texttt{-{}-open-loop true}
flag that disables replanning: ONE plan up front from
INIT\_STATE+TRUTH\_PARAMS, applied for the whole bench. With this
flag the bench produces the smooth recovery$\to$ramp $\Phi$ that
the ground-truth driver does ($\Phi: 0.20 \to 0.12 \to 0.78$ over
14\,d), and $\bar A_{\rm MPC}$ within MC noise of the ground
truth's value ($0.085$ vs $0.091$).

The closed-loop schedule oscillation under the corrected plant call
is a residual MPC-policy concern (each replan can find a different
local optimum), not a kernel/plant/orchestration bug. It is out of
scope for the bug-hunt; the canonical reference for this codebase
is now the open-loop result.

\subsection{Hypothesis 11 (open-loop matches; closed-loop oscillates): replanning is the residual issue}

With \S2.9 (fixed planning horizon) and \S2.10 (\path{advance_subdaily!})
both applied, two regimes can be tested side by side:

\paragraph{Open-loop bench (\texttt{-{}-open-loop true}).}
A new CLI flag in \path{tools/bench_smc_full_mpc_fsa_gpu.jl} disables
replanning: the controller computes ONE plan up front from
INIT\_STATE+TRUTH\_PARAMS, the plant applies it through 27 strides via
\path{advance_subdaily!}, the filter still runs but no replan is
triggered. This mirrors the structure of the ground-truth driver
\path{tools/test_max_A_with_F_barrier.jl} exactly. Result at
$h=15\,$min:
\begin{itemize}[itemsep=2pt]
\item Applied $\Phi$: smooth recovery$\to$ramp $0.20 \to 0.12 \to 0.78$
  over 14\,d, no oscillation. Same qualitative shape as the ground-
  truth driver.
\item $\bar A_{\rm MPC} = 0.085$, $\bar A_{\rm baseline} = 0.082$, $A$
  trajectory dips to ${\sim}0.07$ during recovery and recovers to
  ${\sim}0.10$ by day 14. Within MC noise of the ground-truth driver's
  $\bar A_{\rm MPC} = 0.091$.
\end{itemize}

So the kernel + plant + framework SMC$^2$ + HMC are all working when
exercised by the same single-plan structure as the ground truth. The
canonical reference for this codebase is now this open-loop result.

\paragraph{Closed-loop bench at heavy config.}
The default mode --- replan at end of every $K$ strides, fixed planning
horizon, posterior-mean dynamics params from the filter --- still
produces an oscillating applied $\Phi$ across the bench, with peaks
swinging between $\Phi = 0.5$ and $\Phi = 1.4$. $\bar A_{\rm MPC}$ is
also around $0.087$ (within MC noise of open-loop), so the magnitude
gap to Python is the same; the visual difference is the noise in the
$\Phi$ trace. Two contributions identified:

\begin{enumerate}[itemsep=2pt,leftmargin=*]
\item \textbf{Filter posterior is noisy at $n_{\rm smc} = 32$.} The
  diagnostic prints in \path{bench_smc_full_mpc_fsa_gpu.jl} show
  posterior-mean $\mu_F$ swinging from $0.244$ to $0.261$ (5--7\%
  stride-to-stride) and $\kappa_B$ swinging from $0.0119$ to $0.0130$.
  Each replan plans from a slightly different params vector, which
  can land HMC in a different local minimum of the cost surface.
  Python's reference uses $n_{\rm smc\_particles} = 1024,
  n_{\rm pf\_particles} = 800$ for the filter; Julia's bench uses
  $n_{\rm smc} = 32, K_{\rm per\_chain} = 400$. About $64\times$
  fewer particles.
\item \textbf{Even with TRUTH\_PARAMS forced} (filter posterior bypassed
  for the controller's params), replan-by-replan plans still land in
  different basins because each replan starts from a different
  \texttt{last\_xhat} (filter's smoothed state estimate), and the cost
  surface near the heavy-config MAP appears multimodal (open-loop
  bisection from a low-$F$ state can find plan day-0 $\in [0.7, 1.4]$
  depending on RNG seed). At $n_{\rm smc}=32$ HMC's chains are also
  jittery, which is on top of the filter-state jitter.
\end{enumerate}

\paragraph{Why this is in the closed-loop bench but not the ground-truth driver.}
The ground truth never replans, so the controller commits one plan
in one cost-surface basin. The closed-loop bench replans 13 times in
a 27-stride run; each replan has to find a fresh MAP. The
local-minimum-spread combined with filter-state jitter is what
manifests as the oscillating applied schedule.

\subsection{Performance: Julia under-utilises the RTX 5090 at large $n_{\rm smc}$}

A practical observation while testing fixes for \S 2.11. With
$n_{\rm smc} = 32$ the bench finishes in $\sim 220\,$s; with
$n_{\rm smc} = 1024$ (matching Python's reference filter
configuration) the bench had not finished a single stride after
$\sim 140\,$s of warmup, with GPU utilisation around 51\,\% and only
$\sim 7\,$GB / $32\,$GB of memory used. Linear extrapolation gave a
$\sim 60$--$80\,$min total. Python at the same config takes $\sim
27\,$min, so Julia is roughly $2$--$3\times$ \emph{slower} than Python
at large $n_{\rm smc}$ --- the wrong direction for a pure-GPU port.

The plausible bottlenecks (both visible by reading
\path{run_outer_smc_gpu} in \path{bench_smc_full_mpc_fsa_gpu.jl}
and \path{parallel_hmc_one_move!} in
\path{models/fsa_high_res/gpu_pf.jl}):

\begin{itemize}[itemsep=2pt]
\item Each tempering level (10--11 per stride) calls
  \path{gpu_log_density_batched} once for the cost evaluation and
  then \texttt{num\_mcmc} $\times L$ more times inside the HMC moves
  for the FD gradient batcher (\path{gpu_grads_parallel_chains}
  expands $M$ chains to $M(1+2d) = 61M$ rows, runs the inner-PF on
  all of them, returns to host). Each call has a host-side
  per-chain CPU loop that re-builds the params matrix
  (\texttt{params\_cpu}), an unconditional \texttt{KernelAbstractions.synchronize}
  at every segment of the inner-PF, and a small but per-chain
  \texttt{Array(view(...))} copy back from GPU into a host array
  for the segmented log-likelihood accumulation.
\item Filter resampling and OT rescue paths run a CPU-side per-chain
  branch (line 216, 252 of \path{Filtering/GPUSegmentedPF.jl}).
\item HMC's accept/reject and posterior cloud update are on the host
  in a per-particle loop (\path{parallel_hmc_one_move!}:527--533).
\end{itemize}

None of these are intrinsically GPU-bound limitations; they are
host-side overhead per-chain that scales with $n_{\rm smc}$ even
though the work could be batched. So the practical fix is either
(a) batch these more aggressively, removing the host-side
per-chain loops, or (b) accept the $n_{\rm smc} = 32$ scale and
look for filter stabilisation that doesn't require more particles
(e.g.\ stronger warm-start bridge across windows, or carrying the
controller forward without re-pulling from the filter posterior at
every replan). This is in scope for the next iteration but out of
scope for the bug-hunt that was supposed to be a pure
``why does Python work and Julia not?'' diagnosis.

\subsection{Updated summary table}

\begin{tabular}{@{}p{0.16\textwidth}p{0.5\textwidth}p{0.27\textwidth}@{}}
\toprule
Hypothesis & Test & Outcome \\
\midrule
1.~RBF row-norm    & Removed, re-ran bench         & Real bug. Kept.    \\
2.~Kernel mismatches (NaN guard, $\varepsilon_A$, POST-step) & Aligned to Python, re-ran & Cosmetic. Kept. \\
3.~Bursty/smooth $\Phi$ asymmetry & Routed MPC plant through \texttt{advance!} & No effect. Discarded. \\
4.~Framework controller bugs & 3 synthetic costs with known optima & Framework correct. \\
5.~HMC stuck in wrong basin & Seeded HMC at recovery$\to$overload  & HMC drifts away. Discarded. \\
6.~SDE / decoder / kernel infrastructure & Max-$\int B$ trivial reward & $\Phi \to \Phi_{\max}$ correctly. \\
7.~fp32 cancellation in G1 drift & fp64-promoted substep arithmetic  & Recovery$\to$overload restored at light config. \\
8.~Algebraic rearrangement (Horner $\mu$ + $a_{\rm typ}^{-1}$) & Applied; bit-for-bit + light open-loop + heavy closed-loop & Necessary, not sufficient. Kept. \\
9.~Receding-horizon used shrinking remaining-bench horizon & Replaced with fixed $H_{\rm plan} = T_{\rm bench}$ at every replan & Real bug. Kept. \\
10.~Bench was averaging controller's per-bin $\Phi$ to a daily and re-bursting through plant envelope & Replaced with \path{advance_subdaily!} (controller's smooth $\Phi(t)$ goes straight to plant) & \textbf{Real bug; the residual bottleneck through \S 2.9.} Open-loop now matches ground truth. \\
11.~Closed-loop replan oscillation & Diagnostic prints + open-loop comparison & Open-loop matches ground truth. Closed-loop oscillates due to filter-posterior + cost-surface basin sensitivity. Magnitude gap unchanged. \\
\bottomrule
\end{tabular}

\section{Outstanding}

\begin{itemize}[itemsep=2pt,leftmargin=*]
\item Decide whether to ship the bench as ``open-loop matches ground
  truth, closed-loop is residually noisy'' or to fix the closed-loop
  oscillation. The cheapest closed-loop stabiliser would be carrying
  the controller's posterior forward through a Gaussian bridge across
  replans (\S 7 of the framework reference) so HMC warm-starts in the
  same basin instead of cold-starting from the prior every $K$ strides.
\item Fix the GPU-utilisation regression (\S 2.12). Specifically:
  remove the per-chain host loops in
  \path{run_outer_smc_gpu}, \path{gpu_log_density_batched},
  \path{run_segmented_smc_step!}, and \path{parallel_hmc_one_move!}
  by either pre-allocating a single device buffer for the params /
  accumulators, or moving the per-chain branches inside a single
  \texttt{@kernel} so the GPU stays saturated as $n_{\rm smc}$ grows.
\item Once the controller is producing a stable schedule, switch the
  FD gradient to ForwardDiff dual-number autodiff. Per the bistable
  B3 module's empirical note, ForwardDiff with $d=8$ out-performs
  Enzyme reverse-mode by $\sim 2.7\times$ on a similar PF; the FSA
  controller has the same dimensionality. The seeded-HMC diagnostic
  of \S 2.5 showed gradient quality wasn't the bug, but exact gradients
  buy a $17\times$ throughput improvement per HMC move (one autodiff
  pass vs $1+2d=17$ FD evaluations).
\end{itemize}

\section{What is in the bench today}

\path{version_2_Julia/tools/bench_smc_full_mpc_fsa_gpu.jl}
implements:

\begin{itemize}[itemsep=2pt]
\item \textbf{\S 2.8 algebraic-stable kernel} in
  \path{models/fsa_high_res/gpu_control.jl}: Horner-form $\mu$,
  precomputed $a_{\rm typ}^{-1}$, Kahan-compensated cost accumulators,
  FMA throughout the substep drift. All fp32 except the per-trial
  cost output. Bit-for-bit-matches Python at probe $\theta$ values
  to $3\times 10^{-6}$ relative.
\item \textbf{\S 2.9 fixed planning horizon}: every replan plans
  \texttt{T-days} ahead, regardless of how much bench is left. Mirrors
  Python's \path{plan_horizon_days=T_total_days}.
\item \textbf{\S 2.10 \path{advance_subdaily!} plant call}: the
  controller's per-bin $\Phi$ goes straight to the plant via
  \path{advance_subdaily!} without any daily-averaging or
  burst-envelope re-application. Plant integrates the SDE bin-by-bin
  under the exact signal the cost kernel optimised against.
\item \textbf{\S 2.9 end-of-stride replan strategy}: apply baseline
  $\Phi=1.0$ for the first $K$ strides, then replan every $K$
  strides at end-of-stride. Mirrors Python's
  \texttt{if (s+1) \% K == 0}.
\item \textbf{\texttt{-{}-open-loop true} flag}: disables replanning,
  applies a single up-front plan from INIT\_STATE+TRUTH\_PARAMS for
  the whole bench. Reproduces the ground-truth driver structurally.
\item \textbf{Auto-plot hook} (with \path{Base.invokelatest}
  world-age workaround) that regenerates
  \path{E5_full_mpc_T*_traces.png} and the 30-panel param-trace plot
  on every bench finish.
\end{itemize}

The closed-loop bench at default config (\path{--replan-K 2
--step-minutes 60}) reproduces the qualitative U-shape of Python's
applied $\Phi$ but with an oscillating overlay; \texttt{-{}-open-loop
true} reproduces a smooth recovery$\to$ramp that matches the
ground-truth driver within MC noise.

\section{Running experiments}

This section is for a new agent / user picking up the codebase. Every
experiment in this writeup was produced by one Julia driver:

\begin{center}
\path{version_2_Julia/tools/bench_smc_full_mpc_fsa_gpu.jl}
\end{center}

\subsection{Prerequisites}

\begin{itemize}[itemsep=2pt,leftmargin=*]
\item \textbf{Working directory.} All commands assume you are in
  \path{version_2_Julia/} (so that \path{Project.toml} and
  \path{tools/} are at the right relative paths).
\item \textbf{Julia 1.11+} with the project's environment activated:
\begin{verbatim}
cd version_2_Julia
julia --project=. -e "using Pkg; Pkg.instantiate()"
\end{verbatim}
\item \textbf{NVIDIA GPU} (codebase tested on RTX 5090). The
  KernelAbstractions kernels target CUDABackend.
\item \textbf{No conda env / Python} required. The bench is pure
  Julia. Comparisons against Python use saved \path{.npz} outputs in
  \path{version_2/outputs/...}.
\end{itemize}

\subsection{The minimal invocation}

\begin{verbatim}
julia --project=. tools/bench_smc_full_mpc_fsa_gpu.jl
\end{verbatim}

Defaults: $T = 14\,$d, $h = 15\,$min, replan every stride
($K = 1$), $n_{\rm smc} = 32$, $K_{\rm per\_chain} = 400$, $\text{num\_mcmc} = 5$,
hmc step size $0.025$, hmc leapfrog $8$, seed $42$. Total wall
$\approx 700\,$s on the RTX 5090. Output goes to
\path{outputs/fsa_high_res/g4_runs/T14d_replanK1_h15min_no_infoaware/}.

The auto-plot hook (\S 4 above) writes
\path{E5_full_mpc_T14d_traces.png} (4-panel B/F/A/$\Phi$ trace) and
\path{E5_full_mpc_T14d_param_traces.png} (30-panel filter posterior
trace) into the same output dir before the bench exits.

\subsection{Reproducing Python's reference run (matched config)}

Python's reference plot
\path{version_2/outputs/fsa_high_res/g4_runs/T14d_replanK2_h60min_no_infoaware/}
was produced at $h=60\,$min, $K=2$. To match Python's outer config:

\begin{verbatim}
julia --project=. tools/bench_smc_full_mpc_fsa_gpu.jl \
  --replan-K 2 --step-minutes 60
\end{verbatim}

Total wall $\approx 200\,$s.

\subsection{Reproducing the ground-truth driver structurally}

If you want the smooth recovery$\to$ramp $\Phi$ trace that matches
\path{tools/test_max_A_with_F_barrier.jl} (\S 2.11 above), use
open-loop mode:

\begin{verbatim}
julia --project=. tools/bench_smc_full_mpc_fsa_gpu.jl \
  --open-loop true --step-minutes 15
\end{verbatim}

Total wall $\approx 260\,$s. Disables replanning entirely; one plan
up front from INIT\_STATE+TRUTH\_PARAMS is committed for the whole
bench.

\subsection{All CLI flags}

Every flag has the form \texttt{-{}-name <value>}. Values are parsed
according to the type of the default:

\begin{tabular}{@{}p{0.21\textwidth}p{0.16\textwidth}p{0.55\textwidth}@{}}
\toprule
Flag & Default & Meaning \\
\midrule
\texttt{-{}-T-days} & \texttt{14}~(Int) & Bench duration. Plant rolls
  forward this many days under the controller's schedule. The
  planning horizon at every replan is also \texttt{T-days} (\S 2.9). \\
\texttt{-{}-step-minutes} & \texttt{15}~(Int) & Outer time step in
  minutes. Must divide $1440$. Determines $\text{BINS\_PER\_DAY} =
  1440 / \text{step-minutes}$. Read at module load via
  \texttt{ENV["FSA\_STEP\_MINUTES"]}, so this must be parsed
  \emph{before} the model imports (the driver does this). Common
  values: \texttt{15} (high-res, default), \texttt{60} (matches Python). \\
\texttt{-{}-replan-K} & \texttt{1}~(Int) & Replan every $K$ strides at
  end-of-stride. The first $K$ strides apply baseline $\Phi=1.0$
  (mirrors Python's
  \texttt{daily\_phi\_plan = np.full(T\_total\_days, baseline)}
  initialiser). Then replan happens at \texttt{(s \% K) == 0}. Common
  values: \texttt{1} (replan every stride), \texttt{2} (Python's
  default). Ignored when \texttt{-{}-open-loop true}. \\
\texttt{-{}-N-smc} & \texttt{32}~(Int) & Number of outer-SMC$^2$
  particles for the FILTER. Each particle spawns its own inner-PF.
  Wall-time scales roughly linearly with this value at the current
  Julia implementation (see \S 2.12 — there's a known performance
  regression with per-chain host loops). Python's reference uses
  \texttt{1024}. \\
\texttt{-{}-K-per-chain} & \texttt{400}~(Int) & Inner-PF particles per
  filter chain. Total filter particles = \texttt{N-smc} $\times$
  \texttt{K-per-chain}. Python's reference uses \texttt{800}. \\
\texttt{-{}-num-mcmc} & \texttt{5}~(Int) & Number of HMC moves per
  tempering level inside the FILTER's outer SMC$^2$. Increase if
  HMC acceptance is too high (chain not mixing). The CONTROLLER's
  num\_mcmc is hard-coded to \texttt{ctrl\_num\_mcmc = 10} in the
  bench (line 322). \\
\texttt{-{}-hmc-step-size} & \texttt{0.025}~(Float) & FILTER's HMC
  leapfrog step size. The CONTROLLER uses
  \texttt{ctrl\_hmc\_step = 0.2} (hard-coded, line 324) — different
  from the filter because the controller's $\theta_{\rm ctrl}$ has
  much wider scale. \\
\texttt{-{}-hmc-leapfrog} & \texttt{8}~(Int) & FILTER's leapfrog
  trajectory length. CONTROLLER uses
  \texttt{ctrl\_hmc\_leap = 16} (hard-coded). \\
\texttt{-{}-max-lambda-inc} & \texttt{0.10}~(Float) & Cap on
  $\delta\lambda$ per tempering level for the filter's bisection.
  Smaller = finer schedule, more levels, safer. \\
\texttt{-{}-target-ess-frac} & \texttt{0.5}~(Float) & Target
  effective-sample-size fraction for the filter's tempering bisection.
  Lower = larger $\delta\lambda$ per level, fewer levels but
  riskier. \\
\texttt{-{}-max-temp-levels} & \texttt{30}~(Int) & Hard cap on
  filter tempering levels. The bench currently logs all levels per
  stride. \\
\texttt{-{}-seed} & \texttt{42}~(Int) & Master RNG seed. Plants use
  \texttt{seed\_offset = -{}-seed}; controller noise grids use
  \texttt{seed + 100 + s} per stride. \\
\texttt{-{}-output-dir} & \texttt{""}~(String) & Output directory for
  \path{data.jld2}, \path{manifest.json}, \path{experiment_run.md},
  and the auto-generated PNG plots. If empty (default), the bench
  builds an auto-named path:
  \path{outputs/fsa_high_res/g4_runs/T<days>d_replanK<K>_h<step>min_no_infoaware/}.
  Pass an explicit directory to keep separate runs side-by-side. \\
\texttt{-{}-open-loop} & \texttt{"false"}~(String) & Set to
  \texttt{true} to disable replanning entirely. The bench computes ONE
  plan up front from INIT\_STATE + TRUTH\_PARAMS over the full
  \texttt{T-days} horizon, then applies it through every stride via
  \path{advance_subdaily!}. The filter still runs (so the param-trace
  PNG is still useful) but its posterior never feeds the controller.
  Reproduces \path{tools/test_max_A_with_F_barrier.jl} structurally. \\
\bottomrule
\end{tabular}

\subsection{Hard-coded controller knobs (not currently CLI-exposed)}

These are inside \path{tools/bench_smc_full_mpc_fsa_gpu.jl} around
line 318--328 and have to be edited in source if you want to change
them:

\begin{itemize}[itemsep=2pt,leftmargin=*]
\item \texttt{ctrl\_n\_smc = 1024} --- controller's outer-SMC$^2$
  particles. Already at Python's reference value.
\item \texttt{ctrl\_n\_inner = 128} --- MC trials per cost evaluation
  (= number of CRN noise realisations averaged for each $\theta$
  sample's cost). Python's reference value.
\item \texttt{ctrl\_n\_anchors = 8} --- number of Gaussian RBF anchors
  for the schedule decoder. Determines $\theta_{\rm ctrl} \in
  \mathbb{R}^8$.
\item \texttt{ctrl\_target\_ess\_frac = 0.5}
\item \texttt{ctrl\_num\_mcmc = 10}
\item \texttt{ctrl\_max\_lambda\_inc = 0.20}
\item \texttt{ctrl\_hmc\_step = 0.2}
\item \texttt{ctrl\_hmc\_leap = 16}
\item \texttt{ctrl\_max\_levels = 30}
\item \texttt{ctrl\_target\_nats = 8.0} --- target nat budget for the
  controller's auto-calibrated $\beta_{\max}$.
\item \texttt{ctrl\_sigma\_prior = 1.5} --- prior std on $\theta_{\rm ctrl}$.
\end{itemize}

\subsection{What lands on disk after a run}

In the output directory:

\begin{itemize}[itemsep=2pt,leftmargin=*]
\item \texttt{data.jld2} --- raw data the plot scripts consume:
  \texttt{trajectory\_mpc}, \texttt{trajectory\_baseline} (both
  $(n_{\rm bins}, 3)$ for $B,F,A$), \texttt{Phi\_per\_bin\_mpc} +
  \texttt{Phi\_per\_bin\_baseline}, \texttt{daily\_phi\_per\_stride},
  \texttt{posterior\_particles} ($(n_{\rm strides}, n_{\rm smc},
  30)$), \texttt{param\_names}, \texttt{truth\_params}, plus run
  metadata.
\item \texttt{E5\_full\_mpc\_T<days>d\_traces.png} --- 4-panel
  diagnostic: $B$, $F$, $A$ trajectories (MPC vs constant-$\Phi$
  baseline) + applied daily $\Phi$ schedule.
\item \texttt{E5\_full\_mpc\_T<days>d\_param\_traces.png} --- 30-panel
  filter posterior trace plot, one panel per static dynamics parameter.
\item \texttt{manifest.json} --- machine-readable run config (seed,
  config knobs, wall time, etc.).
\item \texttt{experiment\_run.md} --- human-readable summary.
\end{itemize}

If the auto-plot fails (e.g.\ KernelAbstractions world-age issue), the
\path{data.jld2} is still written. To re-generate plots from
saved \path{data.jld2}:

\begin{verbatim}
julia --project=. tools/plot_state_traces.jl \
  <out_dir>/data.jld2 <out_dir>/E5_full_mpc_T14d_traces.png

julia --project=. tools/plot_param_traces.jl \
  <out_dir>/data.jld2 <out_dir>/E5_full_mpc_T14d_param_traces.png \
  14 15 48
\end{verbatim}

(arguments to \path{plot_param_traces.jl}: \texttt{T-days},
\texttt{step-minutes}, \texttt{stride-bins}.)

\subsection{Suggested experiments}

\begin{enumerate}[itemsep=2pt,leftmargin=*]
\item \textbf{Sanity check the codebase.}
\begin{verbatim}
julia --project=. tools/test_cost_at_theta0.py    # Python side
cd version_2_Julia
julia --project=. tools/test_cost_at_theta0.jl    # Julia side
\end{verbatim}
  Confirms the cost kernel matches Python fp64 to $3 \times 10^{-6}$
  at probe $\theta$ values.

\item \textbf{Fastest closed-loop bench.}
\begin{verbatim}
julia --project=. tools/bench_smc_full_mpc_fsa_gpu.jl \
  --replan-K 2 --step-minutes 60
\end{verbatim}
  $\sim 200\,$s. Default config that matches Python's outer-loop
  config (replan-K=2, h=60min). Currently the closed-loop applied $\Phi$
  oscillates per \S 2.11.

\item \textbf{Smooth recovery$\to$ramp (open-loop).}
\begin{verbatim}
julia --project=. tools/bench_smc_full_mpc_fsa_gpu.jl \
  --open-loop true --step-minutes 15
\end{verbatim}
  $\sim 260\,$s. Reproduces the ground-truth driver's smooth
  schedule.

\item \textbf{High-res, replan every stride.}
\begin{verbatim}
julia --project=. tools/bench_smc_full_mpc_fsa_gpu.jl   # all defaults
\end{verbatim}
  $\sim 700\,$s. $h=15\,$min, $K=1$.

\item \textbf{Stronger filter (warning: slow).}
\begin{verbatim}
julia --project=. tools/bench_smc_full_mpc_fsa_gpu.jl \
  --replan-K 2 --step-minutes 60 --N-smc 256
\end{verbatim}
  Filter at $n_{\rm smc}=256$ (8$\times$ default). Wall scales
  approximately linearly --- expect $\sim 30\,$min until \S 2.12 is
  fixed.

\item \textbf{Plant-only sanity check} (apply Python's exact daily
  $\Phi$ to Julia's plant; quantifies the plant-design gap):
\begin{verbatim}
julia --project=. tools/test_apply_python_schedule.jl
\end{verbatim}
  $\sim 10\,$s. Confirms Julia's plant under Python's schedule
  produces $\bar A_{\rm MPC} \approx 0.087$ vs Python's $0.122$.
\end{enumerate}

\subsection{Saturating the RTX 5090: which flags to push}

Default config gets the RTX 5090 to about 50\,\% utilisation with
$\sim 7\,$GB of $32\,$GB used --- there is a lot of headroom. \S 2.12
explains why blindly bumping \texttt{-{}-N-smc} doesn't proportionally
speed things up (per-chain host loops). But for a novice agent who
wants to push more work onto the GPU and watch utilisation rise,
here are the flags that actually scale GPU work, in order of impact:

\begin{enumerate}[itemsep=4pt,leftmargin=*]
\item \textbf{\texttt{-{}-N-smc} (filter outer-SMC$^2$ particles).}
  Most direct lever. Each filter particle spawns its own inner-PF run,
  so doubling \texttt{N-smc} doubles the total work the GPU has to
  do per stride. The inner-PF kernel launches are with
  $M \times K$ threads; bigger $M$ = bigger launches = better
  occupancy. Combined with the FD gradient batcher
  ($M(1+2d) = 61M$ chains per HMC step), \texttt{N-smc=128} already
  produces $128 \times 61 \times K = 488\mathrm{k}$-particle launches
  with $K=400$, comfortable for the 5090.\\
  \emph{Try:} \texttt{-{}-N-smc 32} (default), then \texttt{64},
  \texttt{128}, \texttt{256}, \texttt{512}. Each step doubles wall.
  Watch \texttt{nvidia-smi} util go up.

\item \textbf{\texttt{-{}-K-per-chain} (inner-PF particles per chain).}
  Total filter work = \texttt{N-smc} $\times$ \texttt{K-per-chain}.
  Bumping K is sometimes preferable to bumping N because it widens
  the per-launch thread count without amplifying host-side per-chain
  loops as much.\\
  \emph{Try:} \texttt{-{}-K-per-chain 400} (default), then
  \texttt{800} (Python's reference), \texttt{1600}.

\item \textbf{\texttt{-{}-step-minutes} (outer time step).} Smaller
  step = more bins per stride = more SDE work per kernel launch.
  $h = 15\,$min produces 96 bins/day; $h = 5\,$min produces 288
  bins/day (3$\times$ work). The substep loop inside the cost
  kernel and the inner-PF propagate kernel both scale linearly with
  bins.\\
  \emph{Try:} \texttt{-{}-step-minutes 60} (smallest work),
  \texttt{15} (default), \texttt{5} (3$\times$ wall vs default).
  Caveat: only values that evenly divide $1440$ are accepted.

\item \textbf{\texttt{-{}-T-days} (bench duration).} More strides =
  more total benches' worth of work. Doesn't widen parallelism per
  launch, but the controller's planning horizon scales with it
  (\texttt{ctrl\_n\_steps} = \texttt{T-days} $\times$ \texttt{BINS\_PER\_DAY})
  so each \emph{controller} call is bigger.\\
  \emph{Try:} \texttt{-{}-T-days 14} (default),
  \texttt{42} (one Banister chronic time constant — single canonical
  horizon), \texttt{84} (two chronic constants).

\item \textbf{Hard-coded controller knobs} (need editing the bench
  source, lines 318--328): \texttt{ctrl\_n\_inner}, \texttt{ctrl\_n\_smc},
  \texttt{ctrl\_hmc\_leap}, \texttt{ctrl\_num\_mcmc}. These already
  default to Python's heavyweight reference values
  ($n_{\rm inner}=128$, $n_{\rm smc}=1024$, $L=16$, mcmc$=10$), so
  the controller side already saturates fairly well.
  Increasing \texttt{ctrl\_n\_inner} beyond 128 widens each cost
  evaluation (more CRN trials averaged per $\theta$).

\item \textbf{Less impactful for saturation} but increases total
  compute (and wall): \texttt{-{}-num-mcmc} (more HMC moves per
  filter tempering level — sequential, doesn't widen parallelism),
  \texttt{-{}-hmc-leapfrog} (longer leapfrog trajectory — same gradient
  batch, just more iterations), \texttt{-{}-max-temp-levels} (caps
  filter levels, usually doesn't bind anyway).
\end{enumerate}

\paragraph{Recommended novice sweep.} In one terminal start a watcher:

\begin{verbatim}
watch -n 0.5 nvidia-smi --query-gpu=utilization.gpu,memory.used \
                         --format=csv
\end{verbatim}

In another, sweep the filter dimensions while keeping everything
else default ($h=60\,$min, $K=2$):

\begin{verbatim}
for N in 32 64 128 256; do
  for K in 400 800 1600; do
    julia --project=. tools/bench_smc_full_mpc_fsa_gpu.jl \
      --replan-K 2 --step-minutes 60 \
      --N-smc $N --K-per-chain $K \
      --output-dir outputs/sweep/N${N}_K${K}
  done
done
\end{verbatim}

Wall time scales roughly linearly with $N \times K$ at the current
implementation; the sweep above is about $50\,$min total at
$h=60\,$min. Read off util/memory from the watcher to see where
saturation kicks in. The bench's auto-plot writes
\path{E5_full_mpc_T14d_traces.png} per cell of the sweep, so you
can also compare quality of the resulting controller schedule
across $N \times K$.

\paragraph{What to look for.}

\begin{itemize}[itemsep=2pt,leftmargin=*]
\item \texttt{utilization.gpu} should approach $90$--$100\,\%$ at the
  larger configs. If it plateaus at $\sim 50\,\%$ regardless of
  $N \times K$, you have hit the per-chain host-loop ceiling
  diagnosed in \S 2.12 --- bumping the flags more won't help; the
  fix is in the source (batch the per-chain loops).
\item \texttt{memory.used} growing predictably with
  $N \times K \times d$ confirms the GPU is allocating
  per-particle buffers; if it does \emph{not} grow with $N$, the
  bench is reusing buffers and the bottleneck is launch overhead,
  not capacity.
\item Per-stride wall-time printed in the bench output
  (\texttt{[stride X/27] Y.Ys ...}) is the practical metric. At
  saturation, stride wall should grow sub-linearly with $N \times K$
  (per-launch fixed costs amortise).
\end{itemize}

\paragraph{Don't bother trying these for saturation.}

\begin{itemize}[itemsep=2pt,leftmargin=*]
\item \texttt{-{}-replan-K}: only changes how often the controller
  fires, not how big each fire is.
\item \texttt{-{}-target-ess-frac}, \texttt{-{}-max-lambda-inc}: tune the
  bisection schedule, change the number of tempering levels but not
  the work per level.
\item \texttt{-{}-seed}: changes RNG; same compute.
\item \texttt{-{}-open-loop true}: skips replanning so total compute
  drops. Use this for the smooth-schedule run, not for saturation
  testing.
\end{itemize}

\subsection{Known gotchas}

\begin{itemize}[itemsep=2pt,leftmargin=*]
\item \textbf{Working directory.} The driver \texttt{include}s
  \path{models/fsa_high_res/FSAHighRes.jl} and the plot scripts via
  paths relative to its own \texttt{@\_\_DIR\_\_}. You must run
  Julia from \path{version_2_Julia/} (or use an absolute path). If
  you run from the repo root, you'll get a
  \texttt{SystemError: opening file ".../tools/..."} immediately.
\item \textbf{Auto-plot world-age.} The bench's auto-plot block uses
  \path{Base.invokelatest} to dispatch to the included plot
  functions. If a future agent edits the plot files and adds new
  exports, those exports won't be visible to the bench's main()
  unless the \path{invokelatest} pattern is preserved. See bench
  lines 627--660.
\item \textbf{One Julia process per bench.} KernelAbstractions
  caches compiled kernels per Julia process. Running multiple
  benches in parallel via \texttt{\&} works but the second one
  pays its own JIT cost. For a sweep, prefer one Julia process
  invoking \path{main(args::Dict)} multiple times in a loop.
\item \textbf{No comparison plot generator.} Comparing a Julia
  output against the Python reference requires running
  \path{tools/compare_to_python.jl} (which reads both \texttt{.npz}
  and \texttt{.jld2}). The bench itself doesn't do this.
\end{itemize}

\end{document}
